% ============================================================
% Chapter 7 — AI Pipeline, Backend Architecture & Machine Learning Design
% Source: PSYCON Engineering Design Document, Vol. 1, Ch. 7
% ============================================================
\section{AI Pipeline, Backend Architecture \& Machine Learning Design}
\label{sec:ai-pipeline-design}

\subsection{AI System Overview}
\label{sec:ai-system-overview}

PSYCON consists of two synchronized AI pipelines whose feature sets are fused into a final inference engine.

\textbf{Physiological AI Pipeline.} Heart rate; heart rate variability (HRV); blood-oxygen trends (SpO\textsubscript{2}, if used); galvanic skin response (GSR/EDA); body temperature; motion/activity.

\textbf{Speech AI Pipeline.} Voice recording; speech characteristics; language features; conversation analysis; acoustic features.

\subsection{Overall Architecture}
\label{sec:ai-overall-architecture}

\Cref{fig:ai-overall-architecture} shows how physiological features from the Wrist Module and speech features from the Audio Module converge in the AI fusion engine.

\begin{figure}[H]
\centering
\begin{tikzpicture}[
    font=\small,
    stage/.style={draw, rounded corners=3pt, minimum width=3.2cm, minimum height=0.75cm, align=center, line width=0.6pt},
    mod/.style={stage, fill=blue!8},
    feat/.style={stage, fill=gray!8},
    fuse/.style={stage, fill=violet!10, minimum width=3.2cm, minimum height=1.2cm},
    arrow/.style={-{Stealth[length=2.2mm]}, line width=0.6pt}
]
\node[mod] (wrist) at (0,4.6) {Wrist Module};
\node[feat, below=8mm of wrist] (physfeat) {Physiological Features};

\node[mod] (audio) at (0,0) {Audio Module};
\node[feat, above=8mm of audio] (speechfeat) {Speech Features};

\node[fuse] (fusion) at (4.6,2.3) {AI Fusion Engine};

\draw[arrow] (wrist.south) -- (physfeat.north);
\draw[arrow] (audio.north) -- (speechfeat.south);
\draw[arrow] (physfeat.east) -| (fusion.north);
\draw[arrow] (speechfeat.east) -| (fusion.south);
\end{tikzpicture}
\caption{Overall AI architecture: physiological features from the Wrist Module and speech features from the Audio Module converge in the AI fusion engine.}
\label{fig:ai-overall-architecture}
\end{figure}

\subsection{Data Acquisition}
\label{sec:data-acquisition-ai}

\textbf{Physiological.} Heart rate, HRV, GSR, temperature, and motion are collected continuously. Every sample receives a timestamp, sensor ID, quality flag, and battery status.

\textbf{Speech.} The target system collects consented speech continuously in an approved study. The current firmware has not yet demonstrated continuous capture; \Cref{fig:speech-pipeline} shows the intended acquisition pipeline.

\begin{figure}[H]
\centering
\begin{tikzpicture}[
    font=\small,
    node distance=0mm,
    field/.style={draw, rounded corners=2pt, fill=orange!8, minimum width=4.2cm, minimum height=0.65cm, align=center, line width=0.5pt}
]
\node[field] (mic) {Microphone};
\node[field, below=5mm of mic] (pcm) {PCM Audio};
\node[field, below=5mm of pcm] (buf) {Buffer};
\node[field, below=5mm of buf] (feat) {Feature Extraction};
\node[field, below=5mm of feat] (inf) {Inference};

\foreach \a/\b in {mic/pcm, pcm/buf, buf/feat, feat/inf}
  \draw[-{Stealth[length=2mm]}, line width=0.6pt] (\a.south) -- (\b.north);
\end{tikzpicture}
\caption{Speech acquisition pipeline: microphone $\rightarrow$ PCM audio $\rightarrow$ buffer $\rightarrow$ feature extraction $\rightarrow$ inference.}
\label{fig:speech-pipeline}
\end{figure}

\subsection{Feature Engineering}
\label{sec:feature-engineering}

\Cref{tab:feature-engineering} summarizes candidate features under consideration for each signal modality.

\begin{table*}[t]
\centering
\caption{Candidate features by signal modality.}
\label{tab:feature-engineering}
\small
\begin{tabularx}{\textwidth}{@{}l Y@{}}
\toprule
\textbf{Modality} & \textbf{Candidate Features} \\
\midrule
Heart & Mean HR; HR variability; resting trend; peak frequency; recovery trend \\
Motion & Activity level; tremor; stability; walking; rest periods \\
Temperature & Mean; drift; sudden changes; long-term baseline \\
GSR & Baseline conductance; slow tonic changes; rapid phasic responses; peak count; recovery slope \\
Speech (acoustic) & Speaking rate; pause duration; pitch statistics; energy/loudness; voice stability; spectral features (e.g., MFCCs) \\
Speech (language) & Vocabulary diversity; sentence length; sentiment; emotion-related language; topic transitions \\
\bottomrule
\end{tabularx}
\end{table*}

\begin{researchnote}
The feature lists above represent candidate features under consideration. Final feature selection should be informed by empirical validation on collected data.
\end{researchnote}

\subsection{Implemented Audio Engineering Baseline}
\label{sec:implemented-audio-baseline}

The software implements one deterministic mono PCM16 extractor named \texttt{psycon\_audio}. Its ordered feature vector contains duration, RMS and peak dBFS, clipping fraction, DC offset, zero-crossing rate, spectral centroid, 85\% spectral rolloff, spectral flatness, mean/standard-deviation/minimum/maximum autocorrelation pitch from 70--400~Hz, relative pitch jitter, bounded voice stability, energy-active frame fraction, pause count, and pause duration. It has no downloaded-model or third-party feature-extractor dependency. Feature changes must update the contract and equivalence tests together.

Every extracted window retains the session and device identifiers, packet sequence, device timestamp, first sample index, sample count, nominal sample rate, SHA-256 source-packet digest, and extractor identity. Before inference, the pipeline returns one explicit state: usable, insufficient audio, clipped, too noisy, corrupt, or missing audio. Missing input is never replaced with zero-valued audio.

Deterministic fixtures cover silence, impulse, a 220~Hz tone, clipped tone, amplitude-modulated speech-like input, and seeded white noise. The local demo passes each fixture through the Protocol v2 decoder and writes a JSON decision report. These tests verify the software contract and abstention behavior only; final thresholds require evidence from the installed INMP441, and continuous capture, light alignment, clock behavior, and transport-loss recovery remain physical integration gates.

The consent-gated \texttt{psycon\_transcription} stage transcribes contiguous acoustic windows marked usable with a locally cached faster-whisper model; rejected windows never silently enter speech-to-text processing. Results preserve timestamped transcript segments, source-sample boundaries, detected language, approximate confidence, engine/model identity, and explicit complete, no-speech, skipped-quality, unavailable, not-requested, or failed states. Recordings are not sent to a third-party transcription API.

The development profile uses the multilingual \texttt{small} model on CPU with int8 computation, replacing \texttt{tiny} because transcript errors directly contaminate downstream language features. The approved production profile runs multilingual \texttt{turbo} with float16 computation on a CUDA GPU inside the PSYCON backend. The model, device, and compute type are deployment configuration and must be recorded with every result.

Before release, compare \texttt{small}, \texttt{medium}, and \texttt{turbo} on representative consented multilingual recordings. Report per-language word or character error rate, language-detection accuracy, end-to-end latency, peak memory, failure rate, and stability of downstream language features. Freeze the selected model only after this evidence is recorded; a server deployment does not waive latency, concurrency, cost, or privacy constraints.

The \texttt{psycon\_language} baseline implements the English features specified in this chapter: vocabulary diversity, sentence count and length, transparent lexicon-based sentiment, emotion-related word counts, topic-transition distance, speech duration, pause duration, and speaking rate. The transcription layer accepts multilingual speech, but the language-feature extractor explicitly abstains for non-English transcripts until validated language-specific tokenization and lexicons exist. Insufficient text and unavailable transcription also abstain explicitly. These features are intended for controlled research comparison; they are not validated emotion recognition, speaker diarization, or clinical interpretation.

The local webpage accepts permission-confirmed WAV recordings, performs decoding and balanced Protocol v2 quality windows of at most two seconds, optionally runs local transcription, and displays acoustic decisions, transcript segments, language features, and provenance. Its permission checkbox is an engineering safeguard rather than a complete participant-consent system, and the localhost development server is not a production data-collection service.

\subsection{Data Fusion}
\label{sec:data-fusion-ai}

The backend aligns sensor and speech data using timestamps. \Cref{fig:data-fusion} shows the fusion process from sensor and speech features to a combined feature vector and the inference engine.

\begin{figure}[H]
\centering
\begin{tikzpicture}[
    font=\small,
    stage/.style={draw, rounded corners=3pt, minimum width=3.6cm, minimum height=0.7cm, align=center, line width=0.6pt},
    src/.style={stage, fill=blue!8},
    vec/.style={stage, fill=gray!8},
    inf/.style={stage, fill=green!10},
    arrow/.style={-{Stealth[length=2.2mm]}, line width=0.6pt}
]
\node[src] (sf) at (0,2.4) {Sensor Features};
\node[src] (spf) at (0,1.2) {Speech Features};
\node[vec] (fv) at (0,-0.2) {Feature Vector};
\node[inf] (ie) at (0,-1.4) {Inference Engine};

\draw[arrow] (sf.south) -- ++(0,-4mm) -| (fv.north);
\draw[arrow] (spf.south) -- (fv.north);
\draw[arrow] (fv.south) -- (ie.north);
\end{tikzpicture}
\caption{Data fusion: timestamp-aligned sensor and speech features are combined into a single feature vector for the inference engine.}
\label{fig:data-fusion}
\end{figure}

\subsection{Backend Components}
\label{sec:backend-components}

\Cref{fig:backend-pipeline} shows the end-to-end backend pipeline, from device data ingestion to the user-facing dashboard.

\begin{figure*}[t]
\centering
\begin{tikzpicture}[
    font=\scriptsize,
    stage/.style={draw, rounded corners=3pt, text width=1.6cm, minimum height=1.1cm, align=center, line width=0.6pt, fill=teal!8, inner sep=2pt},
    arrow/.style={-{Stealth[length=2mm]}, line width=0.6pt},
    node distance=2.5mm
]
\node[stage] (s1) {ESP32};
\node[stage, right=of s1] (s2) {API};
\node[stage, right=of s2] (s3) {Database};
\node[stage, right=of s3] (s4) {Feature Extraction};
\node[stage, right=of s4] (s5) {ML Model};
\node[stage, right=of s5] (s6) {Results};
\node[stage, right=of s6] (s7) {Dashboard};

\foreach \i/\j in {s1/s2, s2/s3, s3/s4, s4/s5, s5/s6, s6/s7}
  \draw[arrow] (\i.east) -- (\j.west);
\end{tikzpicture}
\caption{End-to-end backend pipeline, from ESP32 device data through the API and database, feature extraction, ML model inference, and results, to the user-facing dashboard.}
\label{fig:backend-pipeline}
\end{figure*}

\subsection{Database Structure}
\label{sec:database-structure}

\Cref{tab:database-structure} lists the fields stored for each session.

\begin{table}[H]
\centering
\caption{Fields stored per session.}
\label{tab:database-structure}
\small
\begin{tabularx}{\columnwidth}{@{}Y@{}}
\toprule
\textbf{Field} \\
\midrule
User ID \\
Session ID \\
Timestamp \\
Sensor data \\
Audio metadata \\
Extracted features \\
Model output \\
Confidence score \\
Logs \\
\bottomrule
\end{tabularx}
\end{table}

\begin{warningbox}
For privacy, avoid storing raw audio longer than necessary unless explicit user consent and a clear research need are established. Server-side transcription requires encrypted transport and storage, authenticated access, audit logging, a documented retention and deletion schedule, and separation of participant identity from recordings and transcripts.
\end{warningbox}

\subsection{Machine Learning Workflow}
\label{sec:ml-workflow}

\Cref{fig:ml-workflow} shows the end-to-end machine learning workflow, from raw data collection through to inference.

\begin{figure}[H]
\centering
\begin{tikzpicture}[
    font=\small,
    node distance=0mm,
    field/.style={draw, rounded corners=2pt, fill=blue!6, minimum width=4.2cm, minimum height=0.55cm, align=center, line width=0.5pt}
]
\node[field] (collect) {Collect Data};
\node[field, below=3.5mm of collect] (clean) {Clean Data};
\node[field, below=3.5mm of clean] (sync) {Synchronize};
\node[field, below=3.5mm of sync] (extract) {Extract Features};
\node[field, below=3.5mm of extract] (norm) {Normalize};
\node[field, below=3.5mm of norm] (train) {Train Model};
\node[field, below=3.5mm of train] (val) {Validate};
\node[field, below=3.5mm of val] (deploy) {Deploy};
\node[field, below=3.5mm of deploy, fill=green!8] (infer) {Inference};

\foreach \a/\b in {collect/clean, clean/sync, sync/extract, extract/norm, norm/train, train/val, val/deploy, deploy/infer}
  \draw[-{Stealth[length=2mm]}, line width=0.6pt] (\a.south) -- (\b.north);
\end{tikzpicture}
\caption{End-to-end machine learning workflow, from data collection and cleaning through synchronization, feature extraction, normalization, training, validation, and deployment, to inference.}
\label{fig:ml-workflow}
\end{figure}

\subsection{Model Development}
\label{sec:model-development}

\Cref{fig:model-development} shows the typical model-development workflow.

\begin{figure}[H]
\centering
\begin{tikzpicture}[
    font=\small,
    node distance=0mm,
    field/.style={draw, rounded corners=2pt, fill=orange!8, minimum width=4.2cm, minimum height=0.55cm, align=center, line width=0.5pt}
]
\node[field] (train) {Training Set};
\node[field, below=4mm of train] (valset) {Validation Set};
\node[field, below=4mm of valset] (test) {Testing Set};
\node[field, below=4mm of test] (perf) {Performance Evaluation};
\node[field, below=4mm of perf] (sel) {Model Selection};
\node[field, below=4mm of sel, fill=green!8] (dep) {Deployment};

\foreach \a/\b in {train/valset, valset/test, test/perf, perf/sel, sel/dep}
  \draw[-{Stealth[length=2mm]}, line width=0.6pt] (\a.south) -- (\b.north);
\end{tikzpicture}
\caption{Typical model-development workflow, from training, validation, and testing splits through performance evaluation and model selection to deployment.}
\label{fig:model-development}
\end{figure}

\begin{engnote}
Participant-level separation should be maintained between training and testing sets to reduce data leakage.
\end{engnote}

\subsection{Performance Metrics}
\label{sec:performance-metrics}

Models should be evaluated using metrics appropriate to the task, such as: accuracy; precision; recall; F1 score; ROC-AUC (for classification); and the confusion matrix.

\begin{requirementbox}
Confidence intervals should be reported where possible.
\end{requirementbox}

\subsection{Prediction Output}
\label{sec:prediction-output}

The system should output: a predicted class or score; a confidence estimate; and data-quality indicators.

\begin{warningbox}
Predictions must not be presented as medical diagnoses.
\end{warningbox}

\subsection{Safety \& Clinical Scope}
\label{sec:safety-clinical-scope}

\begin{warningbox}
PSYCON should be described as: a research prototype; a screening/research support tool; \textbf{not} a diagnostic device; and \textbf{not} a replacement for licensed mental health professionals. This distinction is important for research competitions and ethical review.
\end{warningbox}

\subsection{Privacy}
\label{sec:privacy-ch7}

Recommended practices: encrypt transmitted data; use anonymized participant IDs; restrict access to research data; and obtain informed consent before collecting physiological or audio data.

\subsection{Model Updating}
\label{sec:model-updating}

\Cref{fig:model-updating} shows the recommended model-update lifecycle.

\begin{figure}[H]
\centering
\begin{tikzpicture}[
    font=\small,
    node distance=0mm,
    field/.style={draw, rounded corners=2pt, fill=violet!6, minimum width=4.2cm, minimum height=0.6cm, align=center, line width=0.5pt}
]
\node[field] (new) {Collect New Data};
\node[field, below=4mm of new] (review) {Review Quality};
\node[field, below=4mm of review] (retrain) {Retrain};
\node[field, below=4mm of retrain] (validate) {Validate};
\node[field, below=4mm of validate, fill=green!8] (release) {Version Release};

\foreach \a/\b in {new/review, review/retrain, retrain/validate, validate/release}
  \draw[-{Stealth[length=2mm]}, line width=0.6pt] (\a.south) -- (\b.north);
\end{tikzpicture}
\caption{Recommended model-update lifecycle: collect new data, review quality, retrain, validate, and release a new version.}
\label{fig:model-updating}
\end{figure}

\begin{designnote}
Version numbers should be maintained for both datasets and models.
\end{designnote}

\subsection{AI Validation Checklist}
\label{sec:ai-validation-checklist}

\begin{requirementbox}
\textbf{Before deployment, the following must be verified:}
\begin{itemize}
    \item Sensor synchronization verified.
    \item Audio synchronization verified.
    \item Missing-data handling implemented.
    \item Feature extraction validated.
    \item Model evaluated on unseen data.
    \item Confidence outputs verified.
    \item False-positive/false-negative analysis completed.
    \item Documentation completed.
\end{itemize}
\end{requirementbox}

\subsection{Risks}
\label{sec:risks-ch7}

\begin{observationbox}
\textbf{Potential technical risks:}
\begin{itemize}
    \item Motion artifacts affecting physiological signals.
    \item Environmental noise degrading speech quality.
    \item Missing packets.
    \item Sensor contact loss.
    \item Battery-related interruptions.
    \item Dataset imbalance.
    \item Model overfitting.
    \item Domain shift between training and real-world users.
\end{itemize}
\end{observationbox}

\begin{futureworkbox}
How each risk is mitigated should be documented as part of the AI system's risk register.
\end{futureworkbox}

\subsection{Deliverables}
\label{sec:deliverables-ch7}

The final AI package should include: a data collection protocol; a dataset description; feature documentation; the model architecture; the training procedure; the validation methodology; performance metrics; version history; ethical considerations; and user documentation.

\subsection{Chapter Status}
\label{sec:ch7-status}

\begin{validationbox}
\textbf{Completed for this chapter:}
\begin{itemize}
    \item AI architecture
    \item Data pipeline
    \item Feature engineering overview
    \item Backend workflow
    \item Database design
    \item Model development process
    \item Validation strategy
    \item Privacy and ethics
    \item Risk assessment
    \item Deliverables
\end{itemize}
\end{validationbox}
